% !TeX program = xelatex
% Source versions and the comparison cutoff are stated in the manuscript.
\documentclass[UTF8,fontset=fandol,a4paper,11pt]{ctexart}
\usepackage[margin=27mm,headheight=14pt]{geometry}
\usepackage{amsmath,amssymb,amsthm,mathtools,bm}
\usepackage{booktabs,array,longtable}
\usepackage{xcolor,xurl}
\definecolor{linkblue}{RGB}{25,64,117}
\usepackage[colorlinks=true,linkcolor=linkblue,citecolor=linkblue,urlcolor=linkblue,
  bookmarksnumbered=true,pdftitle={A 17/4 Hausdorff bound for six-dimensional Kakeya sets},
  pdfsubject={A derivation from a quantitative four-dimensional input}]{hyperref}
\usepackage{fancyhdr}
\pagestyle{fancy}
\fancyhf{}
\fancyhead[L]{\small 六维 Kakeya 集合的 $17/4$ 下界}
\fancyhead[R]{\small 研究论文整理稿}
\fancyfoot[C]{\thepage}
\renewcommand{\headrulewidth}{0.3pt}
\setlength{\parskip}{0.25em}
\setlength{\emergencystretch}{2em}
\numberwithin{equation}{section}
\theoremstyle{plain}
\newtheorem{theorem}{定理}[section]
\newtheorem{proposition}[theorem]{命题}
\newtheorem{lemma}[theorem]{引理}
\newtheorem{corollary}[theorem]{推论}
\theoremstyle{definition}
\newtheorem{definition}[theorem]{定义}
\newtheorem{inputthm}[theorem]{引用输入}
\theoremstyle{remark}
\newtheorem{remark}[theorem]{注}
\renewcommand{\proofname}{证明}
\newcommand{\R}{\mathbb R}
\newcommand{\T}{\mathcal T}
\newcommand{\Kest}{\mathsf K}
\newcommand{\HH}{\operatorname{dim}_{\mathrm H}}
\newcommand{\ang}{\operatorname{ang}}
\newcommand{\card}[1]{\#\!\left(#1\right)}
\newcommand{\norm}[1]{\left\lVert #1\right\rVert}
\newcommand{\Ldelta}{L_\delta}
\DeclareMathOperator{\mass}{mass}
\title{由四维有限尺度估计导出的\texorpdfstring{\\}{ }六维 Kakeya 集合的 $17/4$ 下界}
\author{}
\date{2026年10月26日}
\begin{document}
\maketitle
\begin{abstract}
以一项明确的四维有限方向质量估计为输入，结合 Katz--Tao 的维数递推，
本文整理得到：任意包含每个方向上一条单位线段的集合 $K\subset\R^6$ 满足
$\HH K\ge17/4$。这里不要求 $K$ 可测、紧或有界，也不要求可测选线。
关键有限步骤是把任意固定环境空间中的角度分布族投影到 $\R^4$，
用高斯小球估计和碰撞图独立集同时保留方向分离、非短投影及每条原线的着色点。
若原平均着色密度为 $\lambda$、方向角度上界常数为 $A$、角度指数为 $s\ge3$，
所得估计保留 $\lambda^2\card{\T}/A$ 和端点对数损失，因而适用于递推所需的非饱和中间族。
随后以外测度和同尺度有限见证将离散估计转回任意集合的 Hausdorff 维数。
本文将四维证明的有限尺度接口与 Katz--Tao 定理作为具名输入，不从四维集合满维数声明直接推导极大函数估计。
相对所核文献中的六维基线 $7-2\sqrt2$，$17/4$ 的指数差约为 $0.078427$。
\end{abstract}
\noindent\textbf{关键词：}Kakeya 集合；Hausdorff 维数；随机投影；有限尺度估计；Katz--Tao 递推。

\medskip
\noindent\textbf{稿件性质与依赖。}
这是既有研究档案 DIM-001 的论文形式整理稿，经过项目内部交叉审查，未经外部同行评审。
四维有限输入的来源、其使用的高级定理及归档证明见附录\ref{app:inputs}。
本文不宣称独立重证整篇四维原稿或 Katz--Tao 递推，不主张绝对首次或最新最优，
也不解决一般维度的满维数猜想。文献比较的核验截止日为2026年10月19日；整理日期不代表完成了新一轮前沿检索。

\section{引言与主要结论}
\label{sec:introduction}

实欧氏空间中的 Kakeya 集合，是指在每个无向方向上都包含一条单位线段的集合。
Kakeya 猜想预言，这种集合在 $\R^n$ 中的 Hausdorff 维数为 $n$。
本文研究的对象始终是实欧氏空间中普通的直线 Kakeya 集合，
不把有限域版本、Minkowski 维数、受限方向结论或附加结构的集合类与之混同。

四维原稿\cite{OAI4}声称四维 Kakeya 集合具有满 Hausdorff 维数。
对高维转移而言，仅有这一集合维数声明尚不足够：递推需要一项能对任意有限线族、
原着色及任意固定环境维数统一调用的定量输入（常数可依赖该环境维数）。
本文使用的是该证明的有限尺度读出及其稀疏方向放大，明确表述为引用输入\ref{input:D4}。
其证明依赖与档案定位在附录\ref{app:inputs}说明。
既有有限证明和成果报告分别见\cite{ArchiveFSD,ArchiveDIM}。

\begin{theorem}[六维下界]\label{thm:main}
采用四维有限输入\ref{input:D4}和 Katz--Tao 递推定理\ref{input:KT}。
若任意集合 $K\subset\R^6$ 在每个欧氏方向上包含一条单位线段，则
\begin{equation}\label{eq:main}
  \HH K\ge\frac{17}{4}=4.25.
\end{equation}
不要求 $K$ 可测、紧或有界，不要求存在可测的方向选线映射。
\end{theorem}

定理的证明分为三个步骤。第一，利用随机线性投影把任意固定 $\R^M$ 中的有限线族送到 $\R^4$。
投影方向可能发生碰撞，投影线段也可能变短；两类损失都必须在使用四维输入之前支付。
第二，保留任意原条数的非饱和估计，得到每个固定实参数 $n\ge4$ 的严格次临界种子
$\Kest(n,4-\zeta)$。对 $0<\zeta<1$，Katz--Tao 定理给出
\begin{equation}\label{eq:roadmap-recurrence}
  \Kest(6,4-\zeta),\quad\Kest(5-\zeta,4-\zeta)
  \quad\Longrightarrow\quad
  \Kest\!\left(6,\frac{17}{4}-\frac{\zeta}{4}\right).
\end{equation}
第三，以同尺度覆盖和外方向测度建立任意集合的 Hausdorff 读出，再令 $\zeta\downarrow0$。

Katz--Tao\cite{KT}的经典六维下界为 $7-2\sqrt2$。
Hickman--Rogers--Zhang\cite[Figure 5]{HRZ}以及 Gao--Liu--Xi\cite[Table 1]{GLX}
的同类型欧氏对照亦明确列出该数值。二者之差为
\begin{equation}\label{eq:comparison}
 \frac{17}{4}-(7-2\sqrt2)
   =2\sqrt2-\frac{11}{4}
   \approx0.078427125>0.
\end{equation}
这一比较只针对上述明确核到的普通六维 Hausdorff 基线，不是对全部当前文献的最优性认证。
本文的贡献归属是四维有限输入的高维后果及所需有限接口的整理；递推本身归于 Katz--Tao。

\paragraph{全文结构。}
第\ref{sec:preliminaries}节固定离散约定与两个具名输入。
第\ref{sec:projection}节证明有限随机投影和原着色保留。
第\ref{sec:seeds}节给出非饱和估计及真实递推参数。
第\ref{sec:hausdorff}节证明任意集合的 Hausdorff 读出并完成定理\ref{thm:main}。
第\ref{sec:scope}节说明范围；附录列出四维输入来源及稀疏化约。

\section{有限尺度约定与引用输入}\label{sec:preliminaries}

\subsection{方向、线族与着色}
所有环境维数 $M$ 都是固定整数，隐含常数可以依赖 $M$ 及固定空间界 $R$，
但不得令 $M$ 随离散尺度变化。令 $0<\delta<1/2$，并记
\[
  \Ldelta=1+\log(1/\delta).
\]
无向单位方向 $u,v$ 的距离取射影角
$\ang(u,v)=\arccos|\langle u,v\rangle|\in[0,\pi/2]$。
方向 $\delta$-分离，是指任意两个不同方向的射影角至少为 $\delta$。

设 $\T$ 是位于固定球 $B(0,R)\subset\R^M$ 中的单位尺度线段或其固定倍 $\delta$ 邻域组成的有限带标签族。
一个着色 $Y$ 为每条原线 $T$ 指定一组与它相交或相距 $O(\delta)$ 的 $\delta$ 格子。
记
\[
 m=\card{\T},\qquad E=\bigcup_{T\in\T}Y(T),\qquad
 S=\sum_{T\in\T}\card{Y(T)}.
\]
这里 $\card{E}$ 表示格子个数，不表示 Lebesgue 体积。
每条原线的长度有固定上界，故
\begin{equation}\label{eq:single-line-upper}
  \card{Y(T)}\le C_{M,R}\delta^{-1}.
\end{equation}
将每个格子移到原线上的最近参数，再按长度 $\delta$ 的参数区间删去重数，
只损失固定因子；因此也可等价地用原线上的分离标记点计数。
全文只在这一固定常数意义下切换格子与标记，不改变小幂指数。

对实数 $s\ge3$ 和 $A\ge1$，所需的方向假设为
\begin{equation}\label{eq:angular-cap}
 \card{\{T\in\T:\ang(v_T,u)\le\theta\}}
 \le A\left(\frac\theta\delta\right)^s,
 \qquad \delta\le\theta\le1,
\end{equation}
对所有方向中心 $u$ 成立。
此处只有上界，没有逐点 Ahlfors--David 下界、凸 Wolff 条件或 sticky 假设。
平均着色密度 $\lambda$ 的含义是
\begin{equation}\label{eq:mean-density}
 S\ge\lambda m\delta^{-1},\qquad 0<\lambda\le1.
\end{equation}
平均下界不能直接替代每条线的下界。

\subsection{四维有限方向质量输入}

\begin{inputthm}[四维有限方向质量估计，档案 DIR-003]\label{input:D4}
固定一个有界四维图像坐标片。对每个 $\varepsilon_D>0$，存在
$\eta_D>0$、$\delta_D>0$ 和 $C_D$，使下述陈述对 $0<\delta<\delta_D$ 成立。
设
\[
 \ell_i(t)=(t,b_i+t v_i),\qquad 0\le t\le1,\qquad
 |b_i|+|v_i|\le R_D,
\]
为有限带标签直线族，$J_i$ 是各自的 $\delta$ 时间网格标记，且
$\card{J_i}\ge\delta^{-1+\eta_D}$。
给定任意正权 $\omega_i$，若每个半开边长 $\delta$ 的斜率格子 $Q$ 都满足
$\sum_{v_i\in Q}\omega_i\le\alpha$，则全部标记访问的格子并集 $E_D$ 满足
\begin{equation}\label{eq:D4}
 \sum_i\omega_i\card{J_i}
 \le C_D\alpha\delta^{-\varepsilon_D}\card{E_D}.
\end{equation}
允许重复标签和任意正权重跨度；常数不依赖最小权重或标签数。
允许网格及访问格子的固定倍扩大，常数相应改变。
\end{inputthm}

该输入是四维证明的有限尺度读出及方向放大的结果，参见\cite{ArchiveSource,ArchiveDIR}，
不是仅由\cite[Theorem 1.1]{OAI4}的集合维数句子直接推出。
本文在其单位权重、输出方向分离的情形使用\eqref{eq:D4}。
特别地，输入保留 $\forall\varepsilon_D\,\exists\eta_D$ 的次序，
不假定 $\eta_D$ 对 $\varepsilon_D$ 有线性下界。

\subsection{Katz--Tao 的估计记号}

令 $N=\delta^{-1}$。沿用\cite[Definitions 5.1, 5.2, 6.1]{KT}的离散线、着色及小幂损失约定。
对于实 $n,d>0$，$\Kest(n,d)$ 表示：在任意固定环境 $\R^M$ 中，
饱和的 $(n-1)$ 方向族与饱和着色给出至少 $N^{d-o(1)}$ 个访问格子。
这里方向族满足所有中心、所有 $N^{-1}\le\theta\le1$ 的角度上界
$N^{o(1)}(N\theta)^{n-1}$，族大小为 $N^{n-1+o(1)}$；
着色饱和指总质量为 $N^{1-o(1)}\card{\T}$，不预设逐线饱和。
参数 $n$ 不是环境维数 $M$，也不要求为整数。

上述 $o(1)$ 只表示原文容许的任意小幂费用。
实际应用时先固定目标输出损失，再选择足小的输入容差，最后取 $N$ 足够大。
第\ref{sec:seeds}节将给出这一用途所需的显式 $\forall\varepsilon\,\exists\eta$ 充分形式，
并只调用严格次临界 $\Kest(n,4-\zeta)$，不声称临界输入输出损失有统一线性模量。

\begin{inputthm}[Katz--Tao 递推，\cite{KT}定理6.2]\label{input:KT}
设 $0<d<n$。若 $\Kest(n,d)$ 与 $\Kest(d+1,d')$ 成立，则
\begin{equation}\label{eq:KT}
  \Kest\!\left(n,\frac{2n+1+d'}4\right)
\end{equation}
成立。输入允许实参数和任意固定环境空间。
\end{inputthm}

该递推的第二处使用可能发生在 $\R^{M+1}$ 中，且需要原文式(40)的非饱和版本
\begin{equation}\label{eq:KT-unsaturated}
  \card{E}\ge
  \frac{\card{\T}}{N^{n-1}}N^{d-o(1)}.
\end{equation}
因此只证明一个固定环境中的饱和结论是不够的。
下文直接保留原条数因子，支付这一接口；递推的完整证明仍引用\cite{KT}。

\section{有限随机投影与原着色保留}\label{sec:projection}

\begin{lemma}[方向碰撞与非短投影]\label{lem:projection}
固定 $M\ge4$、$s\ge3$ 和常数 $a>0$。
若 $m$ 个 $\delta$-分离原方向满足\eqref{eq:angular-cap}，则存在
$P:\R^M\to\R^4$，$\norm P\le C_M$，和一个原标签子族 $\T_*$，使
\begin{align}
 |Pv_T|&\ge c_M &&(T\in\T_*),\label{eq:nonshort}\\
 \ang(Pv_T,Pv_{T'})&>a\delta &&(T\ne T'\in\T_*),\label{eq:projected-separated}\\
 \card{\T_*}&\ge
 \frac{c_{M,s,a}m}{A\delta^{3-s}\Ldelta}.\label{eq:projection-count}
\end{align}
\end{lemma}

\begin{proof}
令 $P$ 的矩阵元为独立标准实高斯变量。
取足够大的 $B_M$，使 $\Pr\{\norm P>B_M\}<1/32$。
对固定单位 $u$，$Pu$ 为标准四维高斯向量；故可取足够小的 $c_M>0$，
使 $\Pr\{|Pu|<c_M\}<1/128$。

若 $u,v$ 的射影角为 $\theta$，改变 $v$ 的符号后可写
\[
 Pv=(\cos\theta)Pu+(\sin\theta)H,
\]
其中 $H$ 是独立于 $Pu$ 的标准四维高斯向量。
条件于 $Pu$，$H$ 在 $(Pu)^\perp$ 上的投影是标准三维高斯。
在 $\norm P\le B_M$ 的事件上，投影方向角不超过 $a\delta$ 蕴含
\[
 \big|\operatorname{proj}_{(Pu)^\perp}H\big|
 \le \frac{C_{M,a}\delta}{\sin\theta}.
\]
三维高斯小球上界因而给出
\begin{equation}\label{eq:collision-probability}
 \Pr\{\norm P\le B_M,\ \ang(Pu,Pv)\le a\delta\}
 \le C_{M,a}\min\{1,(\delta/\theta)^3\}.
\end{equation}
这里没有把范数约束当作独立事件；$Pu=0$ 的事件概率为零。

对每个原方向作角环分解 $\theta\asymp2^j\delta$。
由\eqref{eq:angular-cap}，对应的期望碰撞次数至多
\begin{equation}\label{eq:collision-degree}
 C_{M,s,a}A\sum_{0\le j\le C+\log_2(1/\delta)}2^{j(s-3)}
 \le Q:=C_{M,s,a}A\delta^{3-s}\Ldelta.
\end{equation}
射影角大于一个固定常数的部分，用方向空间的固定有限覆盖及
$m\le C_MA\delta^{-s}$ 同样控制。
当 $s=3$ 时确有对数费用；当 $s>3$ 时保留该稍弱的统一表达已足够。

短方向个数的期望小于 $m/128$；在范数有界事件上计数的碰撞图边数期望至多 $CmQ$。
分别使用 Markov 不等式并合并失败概率，可以同时选择一个 $P$，
使其范数有界、非短方向不少于 $3m/4$，且这些方向的碰撞图边数不超过 $CmQ$。
图中顶点为原标签，投影角不超过 $a\delta$ 时连边。
利用独立集下界
\[
 \alpha(G)\ge\frac{|V(G)|^2}{|V(G)|+2|E(G)|}
\]
得到\eqref{eq:projection-count}，并同时满足\eqref{eq:nonshort}与\eqref{eq:projected-separated}。
\end{proof}

\begin{remark}
引理没有假定投影后的方向分布自动继承某个 Frostman 上界。
碰撞删选之后，$\R^4$ 中的方向分离才通过三维射影角空间的体积打包给出所有尺度的角上界。
引理也不提供任何沿投影核的额外维数。
\end{remark}

\begin{lemma}[平均密度、共同坐标片与访问并集]\label{lem:marks}
设原方向 $\delta$-分离且满足\eqref{eq:angular-cap}，并满足
\eqref{eq:single-line-upper}及\eqref{eq:mean-density}。
可以先选择至少 $c_{M,R}\lambda m$ 条原线，再使用引理\ref{lem:projection}，
获得一个四维共同有界图像坐标片中的子族 $\T'$，在网格 $\rho\asymp\delta$ 上满足
\begin{align}
 m':=\card{\T'}&\ge
 \frac{c_{M,R,s}\lambda m}{A\delta^{3-s}\Ldelta},\label{eq:heavy-output-count}\\
 \card{J'_T}&\ge c_{M,R}\lambda\delta^{-1}
       &&(T\in\T'),\label{eq:individual-output-marks}\\
 \card{E'}&\le C_{M,R,s}\card{E}.\label{eq:output-union}
\end{align}
所有输出标记来自各自原标签；延长支撑直线不增加着色。
\end{lemma}

\begin{proof}
先删除 $\card{Y(T)}<(\lambda/2)\delta^{-1}$ 的轻线。
它们的总质量至多为 $\lambda m\delta^{-1}/2\le S/2$。
结合\eqref{eq:single-line-upper}，留下的重线个数至少为 $c_{M,R}\lambda m$；
每条留下的线都有至少 $(\lambda/2)\delta^{-1}$ 个自己的原标记。
方向上界\eqref{eq:angular-cap}对子族仍成立，无需重新归一化。
对重线族应用引理\ref{lem:projection}。

把每个原着色格子移到其原线的最近参数点，并在原线参数上作固定重数删选。
由于 $|Pv_T|\ge c_M$，每条投影线都有一个坐标的方向分量绝对值不小于固定正数。
按最大分量和符号划分成有限个坐标片，再选其中线数最多的一个。
在此片中以该坐标为共同时间，其他坐标的斜率和截距都有固定界。
原参数间隔为 $\delta$ 的标记，在每个长度 $\rho=C\delta$ 的新时间格内只合并固定个，
故每条仍保留\eqref{eq:individual-output-marks}的标记数。

所有标记位于一个固定有界区域。将新时间轴分成固定有限个单位窗，
让每条线选择自身标记最多的窗，再按窗对线数作鸽巢选择。
逐线标记数与线数各只损失固定比例。平移该共同窗到 $[0,1]$，
必要时延长支撑直线，但只保留窗内已有标记。
这一操作也适用于着色集中在线段端点附近的情形。
预先把引理中的 $a$ 取为足够大的固定数，可使最终斜率在 $\rho$ 尺度分离。

最后，范数有界的 $P$ 将每个原 $\delta$ 格子映入直径 $O_M(\delta)$ 的集合，
后者可由固定多个四维 $\rho$ 格子覆盖。
最近线点移动、坐标片转换和时间格重分只增加固定倍的格子。
因而得到\eqref{eq:output-union}，同时得到\eqref{eq:heavy-output-count}。
没有把不同原线的标记拼接到一条新线上。
\end{proof}

\section{非饱和有限估计与六维递推}\label{sec:seeds}

\begin{proposition}[保留原条数的有限估计]\label{prop:finite}
采用输入\ref{input:D4}。固定 $M\ge4$、$R$、实数 $s\ge3$ 及 $\varepsilon_D>0$。
存在 $\eta_0>0$ 和尺度阈值，使当方向分离的 $m$ 条原线满足
\eqref{eq:angular-cap}、\eqref{eq:mean-density}及 $\lambda\ge\delta^{\eta_0}$ 时，
\begin{equation}\label{eq:finite-sharp-budget}
 \card{E}\ge
 c_{M,R,s,\varepsilon_D}
 \frac{\lambda^2m}{A\Ldelta}\,
 \delta^{s-4+\varepsilon_D}.
\end{equation}
不要求 $m$ 达到方向饱和数。
\end{proposition}

\begin{proof}
由引理\ref{lem:marks}取得四维子族。
其输出斜率分离，故每个斜率格子内的单位权总和至多为固定常数。
先由输入\ref{input:D4}的输出损失取得 $\eta_D$，再取
$\eta_0<\eta_D/2$，最后令 $\delta$ 足够小。
相当网格 $\rho\asymp\delta$ 上，每条输出线的标记数满足四维输入的门槛，
其中所有固定常数都由上述严格余量吸收。
输入\eqref{eq:D4}因而给出
\[
 \card{E}\ge c\card{E'}
 \ge c\delta^{\varepsilon_D}\sum_{T\in\T'}\card{J'_T}
 \ge c\delta^{\varepsilon_D}m'\lambda\delta^{-1}.
\]
代入\eqref{eq:heavy-output-count}即得\eqref{eq:finite-sharp-budget}。
两个 $\lambda$ 分别来自重线条数与每条输出线的原标记数，不能省去其中之一。
\end{proof}

\begin{corollary}[显式小损失形式]\label{cor:uniform}
固定整数 $M\ge4$、空间界 $R$ 及实数 $s\ge3$。对每个 $\varepsilon>0$，
存在 $\eta>0$、$C$ 与 $\delta_0>0$，使
\begin{equation}\label{eq:tolerance-assumptions}
 A\le\delta^{-\eta},\qquad \lambda\ge\delta^\eta
\end{equation}
蕴含
\begin{equation}\label{eq:unsaturated-finite}
 \card{E}\ge C^{-1}m\delta^{s-4+\varepsilon}.
\end{equation}
这些常数与 $m$ 无关。
\end{corollary}

\begin{proof}
先取 $\varepsilon_D<\varepsilon/4$，再取
$0<\eta<\min\{\eta_D/2,\varepsilon/16\}$，必要时进一步缩小。
式\eqref{eq:finite-sharp-budget}中 $\lambda^2/A\ge\delta^{3\eta}$。
最后取 $\delta$ 足够小，以剩余的正指数余量吸收 $\Ldelta$、固定网格及常数，得到结论。
这给出的是 $\forall\varepsilon\,\exists\eta$，没有得到额外的线性模量。
\end{proof}

\begin{corollary}[任意固定实方向参数的种子]\label{cor:seed}
在相同具名输入下，对每个固定实数 $n\ge4$ 与 $0<\zeta<1$，
在任意固定环境空间中成立 $\Kest(n,4-\zeta)$。
\end{corollary}

\begin{proof}
取 $s=n-1\ge3$。
给定目标损失 $\tau>0$，先在推论\ref{cor:uniform}中选择充分小的输出损失，
再取输入容差 $\eta$，使它同时小于 $\tau$ 的一个固定小部分。
若还有饱和条数下界 $m\ge\delta^{-s+\eta}$，则
\[
 \card{E}\ge c\delta^{-4+\varepsilon+\eta}
            \ge c\delta^{-4+\tau}.
\]
这是一项显式的统一有限充分估计。
在 Katz--Tao 的任意小幂约定中，先固定 $\zeta$，
把上述输出损失选为 $\zeta$ 的严格小部分，再使底层所有输入费用低于对应门槛，
即得到所需的严格次临界 $\Kest(n,4-\zeta)$。
在中间族未饱和时，\eqref{eq:unsaturated-finite}仍给出
\[
 \card{E}\ge c\,mN^{4-(n-1)-\varepsilon},
\]
保留了\eqref{eq:KT-unsaturated}所需的原条数因子。
环境维数始终固定；若该环境中不存在满足饱和条数的方向族，相应陈述自然为空真。
\end{proof}

\begin{proposition}[真实的六维参数代入]\label{prop:six}
对每个 $0<\zeta<1$，成立
\begin{equation}\label{eq:six-finite}
 \Kest\!\left(6,\frac{17}{4}-\frac{\zeta}{4}\right).
\end{equation}
\end{proposition}

\begin{proof}
推论\ref{cor:seed}分别给出
\[
 \Kest(6,4-\zeta),\qquad \Kest(5-\zeta,4-\zeta).
\]
第二个实方向参数至少为 $4$，其角度指数为 $4-\zeta\ge3$，
因此确实处于同一有限投影估计的适用范围。
在输入\ref{input:KT}中取 $n=6$、$d=d'=4-\zeta$，便得
\[
 \frac{2n+1+d'}4=\frac{13+4-\zeta}{4}
                  =\frac{17}{4}-\frac\zeta4.
\]
两次种子调用和递推均在固定参数下进行；先固定最终输出损失，
再选有限次调用的输入容差，最后取尺度足够小。
这里既未调用未证的临界统一模量，也未假设方向维数可免费单调传递。
\end{proof}

\section{任意集合的 Hausdorff 读出}\label{sec:hausdorff}

\begin{lemma}[有限估计到任意集合]\label{lem:hausdorff}
设整数 $q\ge2$ 的有限估计 $\Kest(q,d)$ 在前述任意小幂约定下成立。
若 $K\subset\R^q$ 在每个方向上含一条单位线段，则 $\HH K\ge d$。
该陈述不要求 $K$ 或方向选线映射可测。
\end{lemma}

\begin{proof}
反设存在 $t$ 使 $\HH K<t<d$。
先取一个近竖直的开方向片，以斜率 $v\in\R^{q-1}$ 参数化。
每条候选单位线段在时间坐标上有固定正长度的投影，故包含某个有理端点时间区间；
其截距也属于某个整数有界盒。
这些时间区间和截距盒的组合只有可数个，
而其可用斜率集合覆盖整个方向片。
由 Lebesgue 外测度的可数次可加性，某一个组合的可用斜率集合 $S_0$
有正有限外测度 $a>0$。
作一次固定仿射归一化，使共同时间区间为 $[0,1]$，所有相关轨迹位于固定有界区域。
此时尚未对全部斜率选择见证线。

由 $\HH K<t$，可用球 $B_i$ 覆盖 $K$，使半径满足
\begin{equation}\label{eq:hausdorff-cover}
 \sum_i r_i^t\le1,\qquad\sup_i r_i\longrightarrow0.
\end{equation}
对归一化后的固定区域作同样的覆盖即可，固定仿射变换只改变常数。
令 $U_k$ 是所有满足 $2^{-k-1}<r_i\le2^{-k}$ 的球的并集。
它含有至多 $C2^{kt}$ 个球。

取 $c_0>0$ 使 $\sum_{k\ge1}c_0k^{-2}<1$。
任一成功见证线段 $\ell([0,1])\subset K$ 都有
\[
 \sum_k\big|\{u\in[0,1]:\ell(u)\in U_k\}\big|\ge1,
\]
故对某个 $k$，它在 $U_k$ 中的时间长度至少为 $c_0k^{-2}$。
定义 $S_k$ 为存在这样一条有界见证线的斜率集合。
不要求 $S_k$ 可测；仍有 $S_0\subset\bigcup_kS_k$。
再次利用外测度次可加性，存在同一个 $k$ 使
\begin{equation}\label{eq:outer-direction}
 |S_k|^*\ge c_1 a k^{-2}.
\end{equation}
其中 $c_1>0$ 是固定常数。
于是同一尺度同时具有外方向质量下界和逐见证线的时间长度下界。

令 $\delta=2^{-k}$，在有界 $S_k$ 中取一个极大的 $C\delta$-分离有限斜率集。
它的 $C\delta$ 球覆盖 $S_k$，由\eqref{eq:outer-direction}得
\begin{equation}\label{eq:finite-witness-count}
 m\ge c a k^{-2}\delta^{-(q-1)}.
\end{equation}
固定的 $C$ 可保证实际单位方向也在 $\delta$ 尺度分离。
现在仅对这有限个斜率逐个选择一条成功见证线。
方向的所有尺度上界由打包得到
\[
 \card{\{T:\ang(v_T,u)\le\theta\}}
 \le C_q(\theta/\delta)^{q-1}.
\]
每条线在 $U_k$ 中有至少 $c\delta^{-1}k^{-2}$ 个非空时间格，
每格取一个真实线点，再移到最近格子中心。
每个原半径至多 $\delta$ 的覆盖球只贡献固定多个访问格子，因而
\begin{equation}\label{eq:cover-upper}
 \card{E_k}\le C\delta^{-t}.
\end{equation}
最近格子中心到原线的距离至多为 $C\delta$，符合离散支撑线约定。
若采用单位弧长线段约定，可再将固定长度图像线段分为固定多个单位尺度窗口，
逐线保留最富的窗口；只损固定比例，不增加任何标记。

由\eqref{eq:finite-witness-count}与逐线标记数，方向饱和及着色饱和只损失多对数。
先固定 $0<\varepsilon<d-t$，再选有限估计对应的输入容差，最后令 $k$ 足够大。
所有 $k^{-2}$ 和固定常数都在任意小幂门槛内，因此有限下界给出
\begin{equation}\label{eq:cover-lower}
 \card{E_k}\ge c\delta^{-d+\varepsilon}.
\end{equation}
当最大覆盖半径趋于零时，选出的 $k$ 也必须趋于无穷；
式\eqref{eq:cover-upper}与\eqref{eq:cover-lower}遂矛盾。
全过程只在开放球并集的直线原像上计算普通时间长度，未对不可测集合 $K$ 积分，
也未使用可测方向选线。
\end{proof}

\begin{proof}[定理\ref{thm:main}的证明]
对每个 $0<\zeta<1$，命题\ref{prop:six}和引理\ref{lem:hausdorff}给出
\[
 \HH K\ge\frac{17}{4}-\frac\zeta4.
\]
该结论对所有这样的固定 $\zeta$ 成立，令 $\zeta\downarrow0$ 即得\eqref{eq:main}。
这里先完成每个固定参数的有限估计及 Hausdorff 返回，再取参数极限，
没有把一个固定尺度的覆盖数下界直接当作 Hausdorff 维数下界。
\end{proof}

\section{适用范围、归属与限制}\label{sec:scope}

本文结论对普通的任意六维 Kakeya 集合成立的含义，是在明列四维有限输入及
Katz--Tao 具名递推下的推论。作为完整核验责任，读者需同时检查四维有限接口的原证明与所引高级输入；
作者署名、档案哈希、数值试验或内部审查标签都不是正确性证明。

需要特别区分以下三个层次：四维集合的满维数定理；
适用于带标签原标记的四维有限方向质量估计；
由该估计与递推得到的高维后果。本文使用第二层作为真实输入。
其输出对每次新的有限族统一成立，不要求这些族来自同一跨尺度集合。
但它有逐线长标记门槛，不能据此声称任意小着色密度的尖锐 Kakeya 极大函数估计。

方向删选保留原角上界，是未归一化的子族上界；若改为归一化概率律，必须支付剩余质量分母。
平均着色下界先经过重线删选，再进入逐线输入，故公式\eqref{eq:finite-sharp-budget}确有两个 $\lambda$。
投影后只有在碰撞图删选之后才有分离方向，不能免费继承投影 Frostman 控制。
任意环境空间和实角参数是递推的实际需要，并非对结论作形式换名。

相对\cite{KT,HRZ,GLX}中明确列出的普通六维 Hausdorff 基线，
式\eqref{eq:comparison}给出固定正指数改善。
四维原稿\cite[Section 11.2]{OAI4}明确记录的是向更高环境投影所给的下界 $4$。
截至前述核验截止日的有界比较没有识别同一 $17/4$ 的明确报告或可具体采用的同类型支配结果，
但有限检索不证明绝对首创、最新最优或不存在隐含推论。

本文不提出新的 Katz--Tao 递推，也不独立占有四维原稿或其高级输入的贡献。
五维满 Hausdorff 维数、一般维度满维数及真实反例问题均未由本文解决。
本稿适合进一步的外部数学审阅；内部交叉审查与外部同行评审严格区分。

\appendix
\section{四维有限接口的来源与可追溯证明}\label{app:inputs}

\subsection{最初的有限尺度读出}
四维方向质量输入的上游是如下证明层级的有限读出，档案称 SRC-001：
固定有界四维图像坐标片，对每个 $\sigma>0$，存在 $\eta_{\rm src}>0$ 及 $\delta_0>0$，
使方向 $\delta$-分离、线数至少 $\delta^{-3+\eta_{\rm src}}$、
逐线标记数至少 $\delta^{-1+\eta_{\rm src}}$ 的有限族满足
\begin{equation}\label{eq:source-finite}
 \card{E}>\delta^{-4+\sigma}
\end{equation}
（固定常数可以在严格指数余量中吸收）。
这项陈述使用\cite{OAI4}的有限初始化和后续全部模型排除，
不是由 Hausdorff 维数结论作抽象紧性推论。
完整拼接与有限读出记录在\cite{ArchiveSource}，本文不在此重印四维证明。

归档核验的责任包括：对截距和斜率独立作尺度 $\delta$ 的平滑；
由方向分离与线数取得亚幂三维斜率密度上界；
把真实标记区间转成选取时长；控制有限配置复杂度；
在保持未归一化原质量的前提下取得正终端尺度；
最后调用原证明的零窄度、正有限窄度与无穷窄度全部分支。
只核对其中的初始归约不足以认证\eqref{eq:source-finite}。

\subsection{从稠密读出到稀疏与权重形式}
为说明输入\ref{input:D4}并非额外假定方向数饱和，回顾\cite{ArchiveDIR}的有限化约。
先考虑 $m$ 个 $\delta$-分离斜率，其 $\delta$ 邻域在三维斜率空间中的体积为
$\asymp m\delta^3$。
在固定有界范围内作 $L\le C\delta^{-3}/m$ 个平移，
用独立随机平移的期望覆盖量，可选得使这些邻域的并集具有固定正体积的一组平移。
具体地，对固定目标斜率盒中的每点，一次平移命中的概率至少为 $c m\delta^3$，
故取上述数量的平移后，该盒的期望覆盖体积有固定正下界。
于是平移后的带标签斜率中可抽出 $\asymp\delta^{-3}$ 个固定尺度分离方向。

每次斜率平移 $v\mapsto v+a$ 对应时空剪切
\[
 (t,y)\mapsto(t,y+ta).
\]
剪切保留每条复制线自己的时间标记，单个复制的访问格子数至多为 $C\card E$，
全部复制并集至多为 $CL\card E$。
适当重分一个固定倍网格后，取原门槛小于 $\eta_{\rm src}$ 的严格小部分，
即可对放大族使用\eqref{eq:source-finite}，得到
\begin{equation}\label{eq:sparse-source}
 \card E\ge c m\delta^{-1+\sigma}.
\end{equation}
所有剪切范围与坐标界固定，未增加小幂几何费用。

对于一般带权族，把斜率格子作固定有限着色，使同色不同格子中的任意代表在固定尺度分离。
每个非空格子任选一个原标签，保留它自己的标记；代表的访问并集包含于原并集。
由\eqref{eq:sparse-source}，非空格子总数至多为
$C\delta^{1-\sigma}\card E$。
另一方面，每个格子的总权至多 $\alpha$，每条线标记数至多 $C\delta^{-1}$，故
\[
 \sum_i\omega_i\card{J_i}
 \le C\delta^{-1}\alpha\cdot
       (\text{非空斜率格子数})
 \le C\alpha\delta^{-\sigma}\card E.
\]
取 $\sigma<\varepsilon_D$，适当缩小输入门槛并吸收固定常数，即得\eqref{eq:D4}。
这个化约不按权重大小分层，因此不引入依赖最小权重或权重跨度的费用。

\subsection{源证明使用的具名高级输入}
四维内部证明的归档拼接\cite{ArchiveSource}以以下精确版本为高级输入：
\begin{enumerate}
 \item Guth--Wang--Zahl\cite[Theorem 1.1]{GWZ}：三维互异管的凸聚类集合估计，
  使用其原平均着色密度与聚类假设；并非无条件把它替换为极大函数定理。
 \item Wang--Zahl\cite[Definition 4.2, Theorem 5.2]{WZ}的受限平面点积输入。
 \item Carbery--Valdimarsson\cite[Theorem 1]{CV}的端点多线性 Kakeya 输入。
 \item Orponen--de Saxc\'e--Shmerkin\cite[Theorem 1.5]{OSS}的单尺度乘法卷积 Fourier 衰减输入。
\end{enumerate}
此外使用固定迭代次数的组合估计、固定格式实代数与半代数分解、面积公式及标准测度工具。
这些外部高级定理的完整证明未在本文从零重证，作者报告的形式化也未在本项目本地重新运行。
四维输入、上述高级结果与后续 Katz--Tao 递推的责任按这一顺序区分，避免依赖循环。

\subsection{固定版本的研究档案}
可追溯基础为研究仓库2026年10月19日的提交
\texttt{774c3614ec3e810c127c875a9f5abb60748be848}。
下表是文中引用的证明与接口定位；档案 ID 只用于定位，不代替数学陈述。

\begin{longtable}{@{}p{0.16\textwidth}p{0.76\textwidth}@{}}
 \toprule
 档案 & 内容与定位 \\
 \midrule
 \endhead
 RED-001 & R0038完整源归约、有限初始化及内部证明拼接；第20节为任意集合的外测度读出。\\
 DIR-003 & R0009稀疏方向剪切放大及权重格子化约。\\
 FSD-001 & R0029有限高斯投影、原标记与共同时间、任意固定环境和非饱和预算。\\
 DIM-001 & R0109六维成果报告；完整输入、参数链与同类型一手比较。\\
 独立审查 & R0029投影与递推接口、R0038源拼接和任意集合、R0080 retained核心及R0109成果范围的项目内部交叉审查。\\
 \bottomrule
\end{longtable}
本文展开的是投影、有限预算和外测度读出；四维源估计及递推仍按明列输入引用。
因此它是依赖透明的论文整理稿，不以省略源证明的篇幅暗示已在本文完全重证所有引用结果。

\begin{thebibliography}{99}
\bibitem{OAI4}
OpenAI.
\emph{Every Four-Dimensional Kakeya Set Has Full Hausdorff Dimension}.
2026年9月24日预印本，175页；固定仓库版本
\href{https://github.com/openai/math/tree/1581d4cbdd69a0faf34b2089614fe2a814316a6a/preprints/Every-four-dimensional-Kakeya-set-has-full-Hausdorff-dimension-September-24-2026}{openai/math，提交1581d4cb}。

\bibitem{KT}
Nets Hawk Katz and Terence Tao.
\emph{New bounds on Kakeya problems}.
Journal d'Analyse Math\'ematique \textbf{87} (2002), 231--263.
本文所用预印本：\href{https://arxiv.org/abs/math/0102135v1}{arXiv:math/0102135v1}，2001；第5--6节，特别是定理6.2。

\bibitem{HRZ}
Jonathan Hickman, Keith M. Rogers and Ruixiang Zhang.
\emph{Improved bounds for the Kakeya maximal conjecture in higher dimensions}.
\href{https://arxiv.org/abs/1908.05589v1}{arXiv:1908.05589v1}，2019；特别是Figure 5。

\bibitem{GLX}
Chuanwei Gao, Diankun Liu and Yakun Xi.
\emph{Curved Kakeya sets and Nikodym problems on manifolds}.
\href{https://arxiv.org/abs/2503.11574v1}{arXiv:2503.11574v1}，2025；特别是Table 1的欧氏对照。

\bibitem{GWZ}
Larry Guth, Hong Wang and Joshua Zahl.
\emph{A streamlined proof of the Kakeya set conjecture in $\R^3$}.
\href{https://arxiv.org/abs/2601.14411v2}{arXiv:2601.14411v2}，2026；定理1.1。

\bibitem{WZ}
Hong Wang and Joshua Zahl.
\emph{Sticky Kakeya sets and the sticky Kakeya conjecture}.
\href{https://arxiv.org/abs/2210.09581v2}{arXiv:2210.09581v2}；定义4.2及定理5.2。

\bibitem{CV}
Anthony Carbery and Stef\'an Ingi Valdimarsson.
\emph{The endpoint multilinear Kakeya theorem via the Borsuk--Ulam theorem}.
\href{https://arxiv.org/abs/1205.6371v1}{arXiv:1205.6371v1}，2012；定理1。

\bibitem{OSS}
Tuomas Orponen, Nicolas de Saxc\'e and Pablo Shmerkin.
\emph{On the Fourier decay of multiplicative convolutions}.
\href{https://arxiv.org/abs/2309.03068v2}{arXiv:2309.03068v2}；定理1.5。

\bibitem{ArchiveSource}
本项目研究档案。
\emph{R0038：任意集合归约、原证明拼接与有限/Hausdorff出口}。
2026年10月17日；固定版本的
\href{https://github.com/lixiang90/kakeya_research/blob/774c3614ec3e810c127c875a9f5abb60748be848/notes/10_source_audit/source_assembly/R0038_reduction_source_assembly_and_finite_readout_full_proof.md}{RED-001完整证明}。
内部交叉审查，非外部同行评审。

\bibitem{ArchiveDIR}
本项目研究档案。
\emph{R0009：稀疏方向放大与加权方向化约}。
固定版本的
\href{https://github.com/lixiang90/kakeya_research/blob/774c3614ec3e810c127c875a9f5abb60748be848/notes/40_multiscale/plank_inputs/R0009_sparse_direction_amplification.md}{稀疏方向证明}及
\href{https://github.com/lixiang90/kakeya_research/blob/774c3614ec3e810c127c875a9f5abb60748be848/notes/40_multiscale/plank_inputs/R0009_weighted_direction_reduction.md}{权重化约}。

\bibitem{ArchiveFSD}
本项目研究档案。
\emph{R0029：任意环境的有限投影、原标记与六维种子}。
固定版本的
\href{https://github.com/lixiang90/kakeya_research/blob/774c3614ec3e810c127c875a9f5abb60748be848/notes/20_dimension_lifting/finite_seed_transfer/R0029_projection_and_prior_audit.md}{有限投影证明}与
\href{https://github.com/lixiang90/kakeya_research/blob/774c3614ec3e810c127c875a9f5abb60748be848/notes/20_dimension_lifting/finite_seed_transfer/R0029_conditional_seed_and_six_dimensional_bound.md}{种子及递推整合}。
R0029最初的紧集措辞已由R0038任意集合读出扩展。

\bibitem{ArchiveDIM}
本项目研究档案。
\emph{DIM-001：任意六维Kakeya集合的Hausdorff下界17/4}。
2026年10月19日；
\href{https://github.com/lixiang90/kakeya_research/blob/774c3614ec3e810c127c875a9f5abb60748be848/results/reviewed/DIM-001_six_dimensional_Hausdorff_research_report.md}{固定版本成果报告及审查定位}。
\end{thebibliography}

\end{document}
